\section{初步实验设置}

\begin{frame}{实验材料与模型}
  \begin{columns}[T,totalwidth=\textwidth]
    \begin{column}{0.48\textwidth}
      \begin{block}{两类社会决策范式}
        \begin{itemize}
          \item \metric{风险决策框架任务}
          \item \metric{心理适应性决策前提任务}
        \end{itemize}
      \end{block}

      \vspace{0.5em}
      \begin{itemize}
        \item 当前配置中，风险决策框架行为测试含 36 个题项，按正/负框架成对组织。
        \item 决策前提行为测试含 15 个题项，聚合为 6 个心理适应性维度。
      \end{itemize}
    \end{column}
    \begin{column}{0.48\textwidth}
      \begin{block}{中国 LLM 初步验证}
        \begin{itemize}
          \item 共 19 个 LLM 版本。
          \item DeepSeek: 4 个版本。
          \item GLM: 2 个版本。
          \item Qwen: 13 个版本。
          \item 包含思考模式与非思考模式变体。
        \end{itemize}
      \end{block}
    \end{column}
  \end{columns}

  \vspace{0.4em}
  \smallnote{每个行为测试默认以 population-level 方式收集 LLM 模拟被试答案；本次测试配置默认 population 为 50。}
\end{frame}

\begin{frame}{风险决策框架任务：同一结果的两种描述}
  \begin{block}{题目结构}
    每个条件有正框架与负框架两个版本。
    \planB{} 是风险方案；指标关注 LLM 选择 \planB{} 的比例。
  \end{block}

  \vspace{0.4em}
  \begin{tabularx}{\textwidth}{lYY}
    \toprule
    条件 & 正框架 P & 负框架 N \\
    \midrule
    600 人感染致命疾病 &
    \planA{}：200 人会被救。 \newline
    \planB{}：$1/3$ 概率 600 人全被救，$2/3$ 概率无人被救。 &
    \planA{}：400 人会死亡。 \newline
    \planB{}：$1/3$ 概率无人死亡，$2/3$ 概率 600 人全部死亡。 \\
    \bottomrule
  \end{tabularx}

  \vspace{0.4em}
  \smallnote{对应题项示例：`5-P-600life` 与 `6-N-600life`。}
\end{frame}

\begin{frame}{决策前提任务：让社会线索发生变化}
  \begin{tabularx}{\textwidth}{lYl}
    \toprule
    指标 & 题项如何形成差分 & 人类基线含义 \\
    \midrule
    FRP & 确定救 2/200 人 vs 随机抽取 2/200 人被救 & 对公平随机化的偏好 \\
    SAM & 年轻/中年参与者前提下，救老人 vs 救年轻人的差异 & 自我年龄相关调制 \\
    KP & 家庭成员群体 vs 匿名人群，在同等规模下比较 & 亲缘关系是否提高风险方案选择 \\
    SC & 6 人 vs 600 人，在同一前提下比较 & 群体规模是否改变保守性 \\
    GPR & 风险方案收益从 $1/3$ 全救提高到 $2/3$ 全救 & 对赌博收益提高的响应 \\
    SBS & 确定方案从救 2 人提高到救 4 人 & 对确定收益提高的响应 \\
    \bottomrule
  \end{tabularx}

  \vspace{0.4em}
  \smallnote{FRP = Fairness Randomization Preference；SAM = Self-Age Modulation；KP = Kinship Premium；SC = Scale Conservatism；GPR = Gamble Payoff Responsiveness；SBS = Sure Benefit Sensitivity。}
\end{frame}
